% Based on arXiv:2608.24054v1, especially sections 2, 3.4, 3.6, 3.7, 4, 5.2 and 6.
\subsection{Note on arithmetic structure of modulus vacua in flux compactifications --- arXiv:2608.24054}

\textbf{Authors:} Yukiya Furuta, Tatsuo Kobayashi, Tomoyasu Kori, Shuhei Miyamoto, Ryusei Nishida, Hajime Otsuka

\paragraph{主な主張}
RR flux のみを導入した一複素構造モジュラス模型の超対称条件を正則二次方程式に帰着し、トーラス模型では虚軸上の真空の $(\mathrm{Im}\,u)^2$ が Farey 数列とその逆数に対応することを示した\cite{Furuta:2026arithmetic}。

\paragraph{新規性}
真空の縮重度と周囲の空白領域（void）の大きさの数値的相関を見いだし、固定した負の判別式を持つ原始二次形式の類群と Gauss 合成則を用いて、真空とその CP 共役の関係を整理した。

\paragraph{位置付け}
flux 真空の分布を数論から理解し、modular flavor におけるモジュラス真空の選別を考えるための研究である。
扱った CP 対称な一モジュラス模型と非退化条件を満たす多モジュラス模型では、CP を自発的に破る孤立した超対称真空が得られないことも示す。
